\clearpage\begingroup
\ArithmeticHead{4}{L}{PHILIPP MAENNCHEN, GAUSS ALS ZAHLENRECHNER.}
\renewcommand{\thefootnote}{\arabic{footnote})}\setcounter{footnote}{0}
\begin{center}I.\\[3pt]{\large\textls[100]{Stufe der Empirie.}}\end{center}
\noindent\hspace*{\parindent}Dass \textsc{Gauss} sich zu Beginn seiner Forschertätigkeit mit umfangreichen
Zahlenrechnungen beschäftigte, ist nicht verwunderlich. Die Mathematiker
der damaligen Zeit waren vielfach Empiriker; sie ermittelten neue Sätze und
Eigenschaften von Zahlen und Funktionen vorwiegend durch Induktion, und
in einem gewissen naiven Entwicklungszustand betrachtete man einen Satz
\par\clearpage\SourceLeaf{560}\ArithmeticHead{5}{R}{I. STUFE DER EMPIRIE.}
\noindent schon als sichergestellt, wenn die gewählten Beispiele auf eine grössere Anzahl von Dezimalstellen stimmten.

Dass auch \textsc{Gauss} diesen Weg beschritt, lehrt uns z. B. das Bruchstück
einer Jugendübung, die er auf dem Deckblatt seines Exemplars von \textsc{Leistes}\footnote{\textsc{Chr. Leiste}, \textit{Die Arithmetik und Algebra}, Wolfenbüttel 1790; im folgenden kurz als \textsc{Leiste} zitiert. Dieses Buch sollte ihm wohl ursprünglich als Vademecum dienen; er hat es sich daher auch mit Schreibpapier durchschiessen lassen. Sein Geist eilte jedoch so rasch über den Horizont dieses Buches hinaus, dass auf den Schreibseiten zumeist Dinge stehen, die mit dem gedruckten Text daneben keinerlei Beziehung haben.}
\textit{Algebra} aufgezeichnet hat:
\[
\renewcommand{\arraystretch}{1.7}
\begin{array}{r@{\quad}l}
[1]&1+\frac{1}{4}+\frac{1}{9}+\frac{1}{16}+\ldots=\frac{\pi^2}{6}\text{\textsuperscript{2)}}\\\cline{2-2}
{}[2]&1+\frac{1}{9}+\frac{1}{25}+\frac{1}{49}+\ldots=\frac{\pi^2}{8}\text{\textsuperscript{3)}}\\
{}[3]&\frac{1}{4}+\frac{1}{16}+\frac{1}{36}+\frac{1}{64}+\ldots=\frac{\pi^2}{24}\\\cline{2-2}
{}[4]&1+\frac{1}{16}+\frac{1}{49}+\frac{1}{100}+\ldots\\
{}[5]&\frac{1}{4}+\frac{1}{25}+\frac{1}{64}+\frac{1}{121}+\ldots\\
{}[6]&\frac{1}{9}+\frac{1}{36}+\frac{1}{81}+\frac{1}{144}+\ldots=\frac{\pi^2}{54}\\\cline{2-2}
{}[7]&1+\frac{1}{25}+\frac{1}{81}+\frac{1}{169}+\ldots=1{,}074833\text{\textsuperscript{4)}}\\
{}[8]&\frac{1}{8}+\frac{1}{36}+\frac{1}{100}+\frac{1}{196}+\ldots=\frac{1}{32}\pi^2\\
{}[9]&\frac{1}{9}+\frac{1}{49}+\frac{1}{121}+\frac{1}{225}+\ldots=0{,}15886\\
{}[10]&\frac{1}{16}+\frac{1}{64}+\frac{1}{144}+\frac{1}{256}+\ldots=\frac{1}{96}\pi^2.\\\cline{2-2}
\end{array}
\]
\setcounter{footnote}{2}\footnotetext{Vergl. \textsc{Euler}, \textit{Introductio in analysin infinitorum}, 1748, t. 1, Cap. 10; \textsc{Leonhardi Euleri}, Opera omnia, ser. I, vol. VIII, S. 181.}
\setcounter{footnote}{3}\footnotetext{Bei \textsc{Euler} a. a. O. bewiesen.}
\setcounter{footnote}{4}\footnotetext{Der Wert 1,074833 wurde von uns durch eine gut konvergierende Reihenentwicklung geprüft und hat sich als richtig erwiesen. \textsc{Gauss} hat also damals schon mit der Behandlung schwach konvergierender Reihen Bescheid gewusst.}
\noindent Diese Jugendübung gehört noch der empirischen Epoche an; Formel [1]
nimmt \textsc{Gauss} zum Ausgangspunkt für die Herleitung neuer Formeln. Und
zwar gelingt dies bei den Formeln [2], [3], [6], [8] und [10] durch Vergleichung
\par\clearpage\SourceLeaf{561}\ArithmeticHead{6}{L}{PHILIPP MAENNCHEN, GAUSS ALS ZAHLENRECHNER.}
\setcounter{footnote}{0}
\noindent mit der Ausgangsformel; dagegen nicht bei den übrigen. Darum ermittelt er bei [7] und [9] die Zahlenwerte auf einige Dezimalstellen und wird
jetzt vermutlich probiert haben, ob ein solcher Wert mit einem aliquoten
Teil einer Potenz von $\pi$ in Übereinstimmung zu bringen ist. — Das Problem
ist unerledigt geblieben, obgleich die Reihe der reziproken Quadrate in der
Abhandlung von 1812\footnote{Werke III, S. 153.} in der Theorie der $\Pi$-Funktion auftaucht; aber man
erkennt hier schon deutlich \textsc{Gauss}’ Geschicklichkeit in der übersichtlichen
Anordnung, die alle seine Arbeiten kennzeichnet.

Diese Gabe der übersichtlichen Darstellung offenbart sich auch in einer
andern Jugendübung (\textsc{Leiste}, S. 47):
\[
\renewcommand{\arraystretch}{1.25}
\begin{array}{l@{\qquad}l}
\sqrt[3]{1{,}024}= &\\
\begin{array}[t]{r}
+1{,}00000\,00000\phantom{000000}\\
80000000\phantom{000000}\\
853333\phantom{0000}\\
2402953\\\hline
1{,}008000853573\phantom{0000}\\
64013657\phantom{0000}\\\hline
1{,}00793683992\phantom{00000}
\end{array}
&\begin{array}[t]{r}
-0{,}0000640000\phantom{000}\\
13653\phantom{0}\\
44
\end{array}
\end{array}
\]
\begin{center}also\end{center}
\[
\sqrt[3]{2}=1{,}2599210499.
\]
\par\noindent\hspace*{\parindent}Die Erklärung ist einfach:
\[
\sqrt[3]{2}=\sqrt[3]{\frac{2\mathpunct{.}8^3}{8^3}}
=\sqrt[3]{\frac{1024}{8^3}}=\frac{10}{8}(1{,}024)^{\frac{1}{3}}.
\]
\noindent Die Rechenvorteile, um mit Hilfe der Binomialreihe Quadrat- und Kubikwurzeln aus ganzen Zahlen auszuziehen, sind wohl vor \textsc{Gauss} gang und gäbe
gewesen; namentlich hat ja in dieser Hinsicht \textsc{Euler}\textsuperscript{1)} ein grosses Geschick
entfaltet. Aber die praktische Art der Darstellung verdient Bewunderung, da
alles so übersichtlich angeordnet ist und die Zahl der Dezimalstellen ohne
Mühe nach Belieben vermehrt werden kann.
\setcounter{footnote}{2}\footnotetext{\textsc{L. Euler}, \textit{De inventione quotcunque mediarum proportionalium citra radicum extractionem}, Novi Comment. acad. Petrop. 14 (1769): I, 1770, S. 188—214; \textit{Comment. arithm.} 1, 1849, S. 401—413; \textsc{Leonhardi Euleri}, Opera omnia, ser. I, vol. VI, S. 240.}
\clearpage\SourceLeaf{562}\ArithmeticHead{7}{R}{I. STUFE DER EMPIRIE.}
\setcounter{footnote}{0}
\noindent\hspace*{\parindent}Um andere uns erhaltene Jugendübungen richtig zu würdigen, sei noch
folgendes vorausgeschickt:

\textsc{Gauss} hat schon in früher Jugend mit zwei Problemen gleichsam gespielt,
die ihn sein ganzes Leben hindurch beschäftigt haben und für seine Entwicklung bedeutungsvoll geworden sind. Diese Probleme sind das arithmetisch-geometrische Mittel und das Fortschreitungs- bezw. Verteilungsgesetz der Primzahlen. Bei der rechnerischen Beschäftigung mit beiden Problemen spielten
die Logarithmen eine wichtige Rolle. Nun waren dem 14jährigen \textsc{Gauss}
1791 von Gönnern einige mathematische Bücher geschenkt worden, darunter
die Logarithmentafel von \textsc{Schulze}\footnote{\textsc{J. C. Schulze}, \textit{Neue und erweiterte Sammlung logarithmischer .\, .\, .\enspace .\, .\, . Tafeln}, I, II, Berlin 1778.}, in der die dekadischen Logarithmen auf
7 Stellen und die natürlichen (hyperbolischen) Logarithmen der Primzahlen
bis 10000 durch die unermüdliche Arbeit des Artillerieoffiziers \textsc{Wolfram} auf
48 Dezimalstellen berechnet waren. Dies war die erste Logarithmentafel, die
\textsc{Gauss} in die Hände bekam; vorher hat er sich vermutlich ohne Logarithmen
oder mit selbstentworfenen Tabellen beholfen. Entwürfe zu Logarithmentabellen beschäftigen ihn auch später, und es wurde ja schon in der Einleitung auf seine selbstbearbeiteten Logarithmentafeln hingewiesen. Bei der
Bearbeitung solcher Logarithmentabellen geht er auch äusserst praktische
Wege, wenn auch schon ganz ähnliche Überlegungen von \textsc{Huygens} vorausgingen, die \textsc{Gauss} aber damals noch unbekannt waren. Ein Stück der Tabelle, die sich Werke II, S. 502 findet, soll als Probe hier wiedergegeben
werden:
\[
\renewcommand{\arraystretch}{1.65}
\begin{array}{l@{\quad}c@{\quad}l}
a,&\frac{81}{8}&\frac{81}{80}\left(\frac{3^4}{2^4\mathpunct{.}5}\right)\\
b,&\frac{41}{4}&\frac{6561}{6560}\left(\frac{3^8}{2^5\mathpunct{.}5\mathpunct{.}41}\right)\\
c,&2&\frac{1025}{1024}\left(\frac{5^2\mathpunct{.}41}{2^{10}}\right)\\
\multicolumn{3}{c}{.\quad .\quad .\quad .\quad .\quad .\quad .\quad .}
\end{array}
\]
\noindent Das Wesentliche ist, wie man sieht, dass Zähler und Nenner um 1 verschieden sind, so dass die Logarithmenreihe stark konvergiert. Dabei sind
solche Zahlen gewählt, dass jeder neue Bruch höchstens einen neuen Primfaktor liefert. In dem obigen Stück der Tabelle liefert $c$ sogar überhaupt
keinen neuen Primfaktor. Bildet man $\frac{a^2}{bc}$, so ergibt sich:
\par\clearpage\SourceLeaf{563}\ArithmeticHead{8}{L}{PHILIPP MAENNCHEN, GAUSS ALS ZAHLENRECHNER.}
\[
\frac{a^2}{bc}=\frac{3^8}{2^8\mathpunct{.}5^2}\mathpunct{.}
\frac{2^5\mathpunct{.}5\mathpunct{.}41}{3^8}\mathpunct{.}
\frac{2^{10}}{5^2\mathpunct{.}41}=\frac{2^{10}}{10^3}.
\]
\noindent Nun kann man bequem den dekadischen Logarithmus von 2 bestimmen:
\[
10\mathpunct{.}\log 2-3=2\log a-\log b-\log c,
\]
\noindent also
\[
\begin{aligned}
\log 2={}&\frac{1}{10}\Bigl\{3+M\Bigl[2\Bigl(\frac{1}{81}-\frac{1}{2\mathpunct{.}81^2}+\frac{1}{3\mathpunct{.}81^3}-\ldots\Bigr)-\\
&\qquad-\Bigl(\frac{1}{6560}-\frac{1}{2\mathpunct{.}6560^2}+\ldots\Bigr)
-\Bigl(\frac{1}{1024}-\frac{1}{2\mathpunct{.}1024^2}+\ldots\Bigr)\Bigr]\Bigr\},
\end{aligned}
\]
\noindent wo $M$ den Modul des dekadischen Logarithmensystems bedeutet.

Hat man erst einmal den Logarithmus von 2, so kann man der Reihe
nach alle folgenden berechnen.

Nachdem \textsc{Gauss} im Besitz der \textsc{Wolfram}schen Logarithmentabelle war,
konnte er sich an mancherlei weitere Aufgaben wagen, die übrigens auch
noch dem Geist seiner Zeit gemäss sind. Es handelt sich darum, die Logarithmen von Grössen zu bestimmen, die sich bei der zahlenmässigen Berechnung von Reihen ergeben haben. Da solche Grössen auf viele Dezimalstellen
ermittelt waren, so stand der Numerus natürlich nicht vollständig in der Tafel.

\textsc{Gauss} verfuhr nun in der ersten Zeit ganz nach der Anleitung, die er
in \textsc{Lamberts} \textit{Zusätzen zu den logarithmischen etc. Tabellen}, Berlin 1770, S. 53
bis 65 vorfand. Zur Erläuterung diene das Beispiel, das sich auf dem Deckblatt von \textsc{Gauss}’ Exemplar der \textsc{Lambert}schen \textit{Tabellen} von \textsc{Gauss}’ Hand aufgezeichnet findet:
\[
\pi=\frac{22}{7}\mathpunct{.}\frac{2484}{2485}\mathpunct{.}\frac{12983008}{12983009}.
\]
\noindent \textsc{Gauss} hat die Rechnung, an dieser Stelle wenigstens, nicht weiter geführt.
Die Fortsetzung hat man sich offenbar so zu denken, dass man setzt:
\[
\begin{aligned}
\ln\pi={}&\ln 22+\ln 2484-\ln 7-\ln 2485\\
&-\Bigl(\frac{1}{129830008}-\frac{1}{2\mathpunct{.}1298308^2}+{-}\ldots\Bigr),
\end{aligned}
\]
\noindent wobei man mit den zwei ersten Gliedern der eingeklammerten Reihe schon
einen hohen Grad von Genauigkeit erzielt.

\textsc{Gauss} ist zu diesem Ausdruck für $\pi$ in der Weise gekommen, dass er
von dem bekannten Näherungswert $\frac{22}{7}$ ausgeht und mit diesem zunächst in
$3{,}14159265\ldots$ dividiert. Den Quotienten subtrahiert er von 1 und
\par\clearpage\SourceLeaf{564}\ArithmeticHead{9}{R}{II. DAS GAUSSSCHE DIVISIONSVERFAHREN.}
\noindent verwandelt den Rest in einen Bruch mit dem Zähler 1. Der Nenner ist rund
2485. Nun verfährt er in gleicher Weise mit $\frac{22}{7}\mathpunct{.}\frac{2484}{2485}$.
\par\clearpage\endgroup
